% Owned proposed fragment B.132. Central integration is root-owned.
\subsection*{A marked free source pair with full induced profinite topology}

\statement{B.132}{Proved}{The actual inner source loop and corrected d-contact generate a free group of rank two with its full profinite topology}
\label{stmt:B132}
Use the actual ordinary group and complete relation family of B.123,
B.127 and B.131:
\[
 \Gamma_0=W\rtimes S,\qquad S=F(m,v),\qquad
 W=F/N,\qquad F=F(\mathcal U)*T,\qquad T=F(f,g,x),
\]
where
\(\mathcal U=(a,B,d,e,\bar a,\bar B,\bar d,\bar e)\).
The subgroup \(N\) contains every signed source translate and every
ambient conjugate of all forty actual ordinary node words.
For the actual source contacts \(s_1,\ldots,s_4\), set
\begin{equation}\label{eq:B132-source}
 D_s=s_2s_1s_3^{-1}s_1^{-1}s_2^{-1},\qquad
 H_D=\langle m,D_s\rangle\leq\Gamma_0.
\end{equation}
Then the homomorphism from the abstract free group on these two
specified generators to \(\Gamma_0\) is injective. Moreover, \(H_D\)
has its full induced profinite topology in \(\Gamma_0\). Consequently
\begin{equation}\label{eq:B132-completed-pair}
 H_D\cong F_2,\qquad
 \widehat H_D\ \lhook\joinrel\longrightarrow\ \widehat\Gamma_0,
 \qquad
 \overline{H_D}\cong\widehat F_2,
\end{equation}
with the isomorphisms carrying the two free generators to the actual
words \(m,D_s\).

The proof supplies an actual index-two free-seven quotient with its
signed \(m\)-action and extends arbitrary finite quotient kernels of
this free pair through that quotient to \(\Gamma_0\).
It does not identify the full five- or six-source presentation or
decide the selected fixed-source equality.

\begin{proof}
\emph{B.132.1. An actual index-two quotient retaining the whole node family.}
Define
\[
 \chi_\Sigma(f)=\chi_\Sigma(g)=1,\qquad
 \chi_\Sigma(x)=\chi_\Sigma(U_i)=0.
\]
Since \(N\subseteq\ker(F\to T)\), this character factors through
\(W\). Put
\[
 F^{(2)}=\ker(F\to\mathbb Z/2),\quad
 W^{(2)}=F^{(2)}/N,\quad
 T^{(2)}=\ker(T\to\mathbb Z/2),
\]
using the reduction of \(\chi_\Sigma\). The superscript denotes
index two, not free rank two.
For the transversal \(\{1,f\}\), the complete free Schreier basis of
\(F^{(2)}\) consists of
\begin{equation}\label{eq:B132-Schreier-basis}
 U_i,\ fU_if^{-1}\quad(U_i\in\mathcal U),\qquad
 t=f^2,\quad u=gf^{-1},\quad w=fg,\quad x,\quad y=fxf^{-1}.
\end{equation}
In particular the last five words are a free basis of \(T^{(2)}\).

Kill the fourteen family basis elements other than \(B\) and
\(\bar d\) at the even coset, retaining all five elements of
\(T^{(2)}\). The resulting free quotient is
\begin{equation}\label{eq:B132-Q7}
 Q_7=F(B,\bar d)*T^{(2)}.
\end{equation}
This quotient kills the \emph{entire} \(N\). Indeed, by the complete
factorization and character calculation of B.131, the only node seeds
with absolute colors \(B,\bar d\), with \(B\) written first, have
retained heads
\[
 (1,g),\qquad (x^{-1},g^{-1}),
\]
and hence \(\chi_\Sigma\)-differences \(-1,+1\). Every signed source
operator preserves the color and adds the same retained character
to the heads of all eight colors: the additions for
\(m,m^{-1},v,v^{-1}\) are \(0,0,-1,+1\). It follows inductively
that the difference remains odd for every source word. A retained
deck conjugation adds the same value to both heads, and a kernel
conjugation changes only the nonlinear connectors.

A core retained by \eqref{eq:B132-Q7} has color \(B\) or \(\bar d\)
and even retained parity. Thus every whole-orbit node loses at least
one core. Its full conjugate is then one, irrespective of its sign
or nonlinear connector. All other color pairs lose a core
immediately. This proves that every defining node commutator, every
signed source translate and every full normal conjugate is killed.
We therefore have an actual epimorphism
\begin{equation}\label{eq:B132-retraction}
 q:W^{(2)}\longrightarrow Q_7.
\end{equation}
The seven retained words in \eqref{eq:B132-Schreier-basis} give a
literal section of this free-group quotient. Thus \eqref{eq:B132-retraction}
is a retraction; no post-filling fiber-injection assumption is used.

\emph{B.132.2. Both actual source operators on the free-seven quotient.}
Both signed \(m\)-operators preserve the killed family normal
closure: the color is unchanged and each retained head has
\(\chi_\Sigma\)-value zero. They induce inverse automorphisms
\(\alpha,\alpha^{-1}\) of \(Q_7\). In the ordered basis
\((B,\bar d,t,u,w,x,y)\), their exact values are
\begin{equation}\label{eq:B132-seven-actions}
\begin{array}{c|l|l}
 z&\alpha(z)&\alpha^{-1}(z)\\\hline
 B&t^{-1}wBw^{-1}t&uBu^{-1}\\
 \bar d&t^{-1}w\bar d w^{-1}t&u\bar d u^{-1}\\
 t&t^{-1}wut&uw\\
 u&w^{-1}t&utw^{-1}u^{-1}\\
 w&t^{-1}wt&uwu^{-1}\\
 x&xw^{-1}t&xu^{-1}\\
 y&t^{-1}wyu^{-1}w^{-1}t&uytw^{-1}u^{-1}.
\end{array}
\end{equation}
These formulas are obtained by substituting the complete actual
free-eleven words of B.123 into the index-two Schreier scan.
For a signed letter, let \(\epsilon\in\{0,1\}\) be its starting
coset. A positive generator \(z\) emits
\(f^\epsilon zf^{-\epsilon'}\), where
\(\epsilon'=\epsilon+\chi_\Sigma(z)\pmod2\); a negative letter
emits the inverse of the positive Schreier generator based at its
ending coset. Delete precisely the fourteen killed basis elements.
This specifies the calculation on the whole source words, including
every nonlinear correction. All twenty-one original basis generators
satisfy both quotient diagrams, and the two columns of
\eqref{eq:B132-seven-actions} compose to the identity on all seven
surviving generators in both orders.

The restriction to \(T^{(2)}\) equals the restriction of the
following automorphism of the larger retained free group:
\begin{equation}\label{eq:B132-retained-alpha}
 \alpha(f)=f^{-1}gf,\qquad \alpha(g)=f,\qquad
 \alpha(x)=xg^{-1}f.
\end{equation}
Its inverse sends \(f\mapsto g\), \(g\mapsto gfg^{-1}\) and
\(x\mapsto xfg^{-1}\). Substitution on the five displayed retained
Schreier generators gives exactly \eqref{eq:B132-seven-actions}.
Consequently there is an actual epimorphism
\begin{equation}\label{eq:B132-mapping-torus-quotient}
 P:=W^{(2)}\rtimes\langle m\rangle
 \longrightarrow
 G_7:=Q_7\rtimes_\alpha\langle m\rangle,
\end{equation}
identity on \(m\). The free section of
\eqref{eq:B132-retraction} is not asserted \(m\)-equivariant;
\eqref{eq:B132-mapping-torus-quotient} is an epimorphism, not an
asserted semidirect-product retraction.

\emph{B.132.3. The exact source word and a simultaneous basis.}
Set
\[
 \kappa=f^{-1}g=t^{-1}w\in T^{(2)},\qquad
 D_0=x\bar d x^{-1}.
\]
The table gives \(\alpha(B)=\kappa B\kappa^{-1}\),
\(\alpha(\bar d)=\kappa\bar d\kappa^{-1}\) and
\(\alpha(x)=x\kappa^{-1}\). Hence \(\alpha(D_0)=D_0\).
Replacing \(\bar d\) by \(D_0\) is a free-basis change, with
inverse \(\bar d=x^{-1}D_0x\). Thus
\begin{equation}\label{eq:B132-D0-basis}
 Q_7=\langle D_0\rangle*R,\qquad
 R=F(B)*T^{(2)},\qquad \alpha(R)=R.
\end{equation}
The same complete raw-source scan gives, in the original contact order,
\begin{equation}\label{eq:B132-contact-images}
 (q(s_1),q(s_2),q(s_3),q(s_4))
   =(B^{-1},1,x\bar d^{-1}x^{-1},1).
\end{equation}
In particular the actual word \eqref{eq:B132-source} has image
\begin{equation}\label{eq:B132-D-image}
 d_*=q(D_s)=B^{-1}D_0B.
\end{equation}
This uses the raw source marking, not a substitution of a plain
\(d\)-norm for the corrected source coordinate.

\emph{B.132.4. Every integer translate belongs to one free Schreier basis.}
Choose the ordered conjugators \(\kappa_n\in T^{(2)}\) for all
\(n\in\mathbb Z\) by
\[
 \alpha^n(B)=\kappa_nB\kappa_n^{-1},\qquad
 \kappa_0=1,\qquad
 \kappa_{n+r}=\alpha^n(\kappa_r)\kappa_n.
\]
Put \(q_0=gf\) and \(q_1=fg=fq_0f^{-1}\). From
\eqref{eq:B132-retained-alpha}, \(\alpha(q_0)=q_0\),
\(\alpha^2|_{F(f,g)}=\operatorname{Inn}(q_0^{-1})\), and
\(\kappa_2=q_0^{-1}q_1\). The recursion and its inverse give
\begin{equation}\label{eq:B132-all-integer}
 \kappa_{2n}=q_0^{-n}q_1^n,
 \qquad
 \kappa_{2n+1}=q_0^{-n}(f^{-1}g)q_1^n
 \quad(n\in\mathbb Z).
\end{equation}

These words are pairwise distinct. Their abelianizations separate
even from odd indices. For even indices rewrite
\(\kappa_{2n}=q_0^{-n}fq_0^nf^{-1}\). Equality at two indices
would put a power \(q_0^r\), with \(r\ne0\), in the centralizer
\(\langle f\rangle\) of the free basis element \(f\), contrary
to its \(g\)-exponent. For odd indices use
\[
 \kappa_{2n+1}=q_0^{-n}(f^{-1}gf)q_0^nf^{-1}.
\]
Equality would make a nonzero power of \(q_0\) centralize
\(f^{-1}gf\). Conjugating by \(f\) would put a nonzero power of
\(q_1\) in \(\langle g\rangle\), contrary to its
\(f\)-exponent. The two basis-element centralizer statements follow
directly by reduced-word normal form in \(F(f,g)\).

Let \(B_n=\kappa_nB\kappa_n^{-1}\in R\).
The \(B_n\) are distinct: an equality would put a nontrivial
element of \(T^{(2)}\) in the centralizer of \(B\) in
\(F(B)*T^{(2)}\), excluded by its free-product normal form.
The kernel of the retraction
\(\langle D_0\rangle*R\to R\) has the simultaneous free basis
\[
 rD_0r^{-1}\qquad(r\in R).
\]
Therefore
\begin{equation}\label{eq:B132-orbit-basis}
 \alpha^n(d_*)=B_n^{-1}D_0B_n\quad(n\in\mathbb Z)
\end{equation}
is a family of distinct members of that one basis, and is freely
independent. The action of \(m\) shifts its index by one.
Its free group semidirect the shift \(\mathbb Z_m\) is the free
group on \(m,d_*\): a word of nonzero \(m\)-exponent is detected
by \(\mathbb Z\), and an exponent-zero word rewrites in exactly
the independent basis \eqref{eq:B132-orbit-basis}.
The composite of the abstract free group on \(m,D_s\) with
\eqref{eq:B132-mapping-torus-quotient} is thus injective. This
proves the actual discrete freeness of \(H_D\).

\emph{B.132.5. Finite quotients with a prescribed orbit divisibility.}
For every integer \(n\geq1\), there is a characteristic
finite-index subgroup \(C\triangleleft R\) such that the orbit
of \(\bar B\) under the induced automorphism of \(\bar R=R/C\)
has length divisible by \(n\).

First take the product \(Q\) of two finite images of \(T^{(2)}\).
The first records the \(f\)-exponent modulo two. Each odd
\(\kappa_{2r+1}\) has \(f\)-exponent \(-1\), so it is
nontrivial in this image. The second is the restriction to
\(T^{(2)}\) of
\[
 F(f,g,x)\longrightarrow\operatorname{UT}_3(\mathbb Z/n),\qquad
 f\longmapsto(1,0,0),\quad g\longmapsto(0,1,0),\quad x\longmapsto1,
\]
where
\((a,b,c)(a',b',c')=(a+a',b+b',c+c'+ab')\).
A direct multiplication gives
\[
 q_0^{-r}fq_0^rf^{-1}\longmapsto(0,0,r).
\]
Thus \(\kappa_{2r}\) can be one in \(Q\) only if \(n\mid r\).
This includes \(n=1\), when the second image is trivial.

Map \(R=F(B)*T^{(2)}\) to the finite wreath product
\(C_2^{Q}\rtimes Q\), sending \(B\) to the nonzero lamp at the
identity and \(T^{(2)}\) to \(Q\). Conjugation moves that one
lamp. Equality of the images of \(B_r\) and \(B\) therefore
implies \(\kappa_r=1\) in \(Q\), and hence \(2n\mid r\).
Let \(K_R\) be the kernel of this finite map. Since \(R\) is
finitely generated, there are finitely many subgroups of index at
most \([R:K_R]\): each is obtained from an action on a finite set
of that bounded size. Their intersection is a characteristic
finite-index subgroup \(C\subseteq K_R\).
The automorphism \(\alpha\) descends to an automorphism
\(\bar\alpha\) of \(\bar R=R/C\). If
\(\bar\alpha^r(\bar B)=\bar B\), equality holds in the previous
finite target, so \(2n\mid r\). Its finite orbit length is
therefore divisible by \(n\). The initial wreath quotient itself
need not be \(\alpha\)-invariant; the characteristic refinement
is the step supplying invariance.

\emph{B.132.6. Extension of an arbitrary finite target of the pair.}
Let \(H_*:=\langle m,d_*\rangle\leq G_7\), and let
\(\psi:H_*\to A\) be any homomorphism to a finite group.
Write
\[
 u_A=\psi(m),\qquad z_A=\psi(d_*),\qquad
 n=\operatorname{ord}(u_A).
\]
Choose \(\bar R,\bar\alpha\) as in B.132.5, with orbit length
\(\ell\) of \(\bar B\) divisible by \(n\).
Use the left regular action of \(\bar R\) on \(A^{\bar R}\):
\[
 (r\cdot F_0)(s)=F_0(r^{-1}s).
\]
On \(A^{\bar R}\rtimes\bar R\), define
\begin{equation}\label{eq:B132-theta}
 \Theta(F_0,r)=(F_\Theta,\bar\alpha(r)),\qquad
 F_\Theta(s)=u_A F_0(\bar\alpha^{-1}(s))u_A^{-1}.
\end{equation}
This is an automorphism, because \(\bar\alpha\) intertwines the
left regular actions. Choose a positive integer \(M\) divisible
by both \(\operatorname{ord}(\bar\alpha)\) and
\(\operatorname{ord}(u_A)\). Then \(\Theta^M=1\), and we have
the finite group
\[
 Z_A=(A^{\bar R}\rtimes\bar R)\rtimes_\Theta C_M,
\]
whose final cyclic generator is denoted \(m_*\).

Define the function \(F_*:\bar R\to A\) by
\begin{equation}\label{eq:B132-lamp}
 F_*(\bar\alpha^j(\bar B))=u_A^jz_Au_A^{-j},
 \qquad F_*=1\ \hbox{off this orbit}.
\end{equation}
It is well-defined because \(n\mid\ell\), and is fixed by
the function action in \eqref{eq:B132-theta}. The assignments
\[
 D_0\longmapsto(F_*,1),\qquad
 R\longmapsto\bar R,\qquad m\longmapsto m_*
\]
therefore define a homomorphism \(\Phi:G_7\to Z_A\): the
free-product map on \(\langle D_0\rangle*R\) intertwines
\(\alpha\) with \(\Theta\) on both factors.

The image of \(H_*\) lies in the subgroup
\(A^{\bar R}\rtimes C_M\) of \(Z_A\), with trivial
\(\bar R\)-component. On this subgroup, and only on this
subgroup, use
\begin{equation}\label{eq:B132-evaluation}
 \operatorname{Ev}(F_0,m_*^j)=F_0(1)u_A^j.
\end{equation}
This is a homomorphism to \(A\). Indeed \(\bar\alpha\) fixes
the identity index, so \(\Theta\) acts on its coordinate by
conjugation by \(u_A\); also \(u_A^M=1\). The left-action
convention gives
\[
 \operatorname{Ev}(\Phi(d_*))
   =(\bar B^{-1}\cdot F_*)(1)=F_*(\bar B)=z_A,
 \qquad \operatorname{Ev}(\Phi(m))=u_A.
\]
Thus \(\operatorname{Ev}\Phi|_{H_*}=\psi\), and
\begin{equation}\label{eq:B132-kernel-containment}
 \ker(\Phi)\cap H_*\subseteq\ker(\psi).
\end{equation}
This proves the full induced topology of \(H_*\) in \(G_7\).
It is a finite-kernel refinement; it does not assert an extension
\(G_7\to A\) of \(\psi\), since \(\operatorname{Ev}\) was
not defined on all of \(Z_A\).

\emph{B.132.7. Transfer to the actual ambient group and completion.}
Pull back \(\Phi\) along \eqref{eq:B132-mapping-torus-quotient}.
Because this quotient restricts to the proved isomorphism
\(H_D\to H_*\), \eqref{eq:B132-kernel-containment} proves
full induced topology of \(H_D\) in \(P\).

We finally check the actual ambient comparison. Put
\(P_{\rm big}=W\rtimes\langle m\rangle\).
For any finite homomorphism \(\rho:P_{\rm big}\to A_0\),
choose a characteristic finite-index subgroup
\(W_0\subseteq\ker(\rho)\cap W\). It exists by the bounded-index
intersection argument, because \(W\) is finitely generated.
The entire actual free source group \(S=F(m,v)\) acts on the
finite group \(W/W_0\). Let \(n_0=\operatorname{ord}(\rho(m))\).
Map \(S\) to the product of its finite action group and
\(C_{n_0}\), giving \(m\) cyclic coordinate one and \(v\)
coordinate zero. This is a homomorphism because the two source
lifts freely generate \(S\).
The resulting finite semidirect quotient of \(\Gamma_0\) has
kernel on \(P_{\rm big}\) contained in \(\ker\rho\): a kernel
element \(wm^j\) has \(w\in W_0\) and \(n_0\mid j\).
Thus \(P_{\rm big}\) has full induced topology in \(\Gamma_0\).

The subgroup \(P=W^{(2)}\rtimes\langle m\rangle\) has index
two in \(P_{\rm big}\). Taking a finite-index core shows that
every finite quotient kernel of \(P\) contains the restriction
of a finite quotient kernel of \(P_{\rm big}\).
Full topology is transitive under these kernel refinements, so
\(H_D\) has its full topology in \(\Gamma_0\). The induced
map of its completion is injective, with image its actual closure.
The discrete free-pair identification then proves
\eqref{eq:B132-completed-pair}.

No residual-finiteness shortcut, equivariant section of the
free-seven retraction, or completion of an infinite-rank abstract
kernel was used. The conclusion concerns these two marked source
words only.
\end{proof}

\begingroup\raggedright\footnotesize
Exact actual quotient and source-word certificate:
\path{work_CE/code/compile_even_sum_m_quotient26.py},
\path{work_CE/code/compile_even_sum_m_quotient26.out},
\path{work_CE/next_symmetric_nodes/even_sum_m_quotient26.json}.
The checks include both twenty-one-generator quotient diagrams,
both free-basis changes, all inverse action compositions and the raw
source images. The complete all-integer independence and arbitrary
finite-target extension are proved above, without a finite orbit cutoff.
Original full proofs and independent reviews are retained in
\path{work_CE/even_sum_m_free_subgroup26.txt},
\path{work_A/actual_free_source_pair_full_topology26.md},
\path{work_A/review_even_sum_m_free_subgroup26.md},
\path{work_DE/next_marginal_source/review_even_sum_m_free_subgroup26.md},
and \path{work_CE/review_free_source_pair_full_topology26.txt}.
\par\endgroup
